\documentclass[addpoints,answers,a4paper,11pt]{exam}
\usepackage{amssymb,amsmath}
\usepackage[usenames]{color}
\usepackage{multicol,fancybox,tabularx,textcomp,tabulary,pbox,ifthen}
\usepackage[dvipsnames]{xcolor}
\usepackage{blkarray}
\usepackage{booktabs}
\usepackage{multirow}
\usepackage{multicol}
\usepackage{makecell}
\usepackage{tikz}

%***************** solutions color etc preferences ********************
\unframedsolutions
%\shadedsolutions
%\definecolor{SolutionColor}{rgb}{.98,.98,.98}
%\renewcommand{\solutiontitle}{\noindent\textsf{Soln.:}\par\noindent}

%******************** Course Details *********************************
\newcommand{\testname}{End-Semester Examination}
\newcommand{\courseno}{MA 324}
\newcommand{\coursename}{Statistical Inference and Multivariate Analysis}
\newcommand{\instiname}{IIT GUWAHATI}
\newcommand{\testdate}{May 8, 2026}
\newcommand{\testtime}{14:00--17:00 IST}
\newcommand{\semester}{2025-26 Even}

%*****************************************************************************
\newcommand{\emptyframedbox}{\framebox{\rule{0pt}{10pt}{\rule{10pt}{0pt}}}}
\newcommand{\emptyrollno}{\emptyframedbox\emptyframedbox\emptyframedbox\emptyframedbox\emptyframedbox\emptyframedbox\emptyframedbox\emptyframedbox\emptyframedbox}
\newcommand{\emptyname}{\emptyframedbox\emptyframedbox\emptyframedbox\emptyframedbox\emptyframedbox\emptyframedbox\emptyframedbox\emptyframedbox\emptyframedbox\emptyframedbox\emptyframedbox\emptyframedbox\emptyframedbox\emptyframedbox\emptyframedbox\emptyframedbox\emptyframedbox\emptyframedbox\emptyframedbox}

%************* FOR FOOTER AND HEADER **********************************
% page foot and header
\pagestyle{headandfoot}

\firstpagefooter{}{}{\tiny Page \thepage\ of \numpages\normalsize}
\runningheader{\semester}{\testname}{\courseno}
\runningheadrule
%\runningfootrule
\runningfooter{}{}{\tiny Page \thepage\ of \numpages\normalsize}
%****************** TO REDUCE MARGINS*******************
\extrawidth{.65in}
\extraheadheight{-0.25in}
\extrafootheight{-0.40in}
%**END--OF--PROOF: stretch is a rubber space which fills. * ensures it works at beginning of line
\newcommand{\eop}{\hspace*{\stretch{1}}{\scriptsize\maltese}}
\renewcommand{\eop}{}
%stretch is a rubber space which fills. * ensures it works at beginning of line
\newcommand{\push}{\hspace*{\stretch{1}}}
%to separate advanced questions
%\newcommand{\advanced}{\fullwidth{\raisebox{-2pt}{$\ll$}{\hrulefill}\raisebox{-2pt}{\tiny\fbox{ADVANCED}}{\hrulefill}\raisebox{-2pt}{$\gg$}}}
\newcommand{\advanced}{\fullwidth{\hrulefill\raisebox{-2pt}{\tiny\fbox{ADVANCED}}\hrulefill}}
%ignore extra spaces around enviroments
%\ignorespaces
\parindent 0pt
%**********some exam package options
%\usehorizontalhalf
%\marginpointname{}
%\marginbonuspointname{}
%\pointsinmargin
%\pointformat{[\themarginpoints $^\mathrm{pnts.}$]}
%\bonuspointformat{[\themarginpoints $^\mathrm{xtra}$]}

% %*********** step mark macro *************************
% \newif\ifstepmark
% \newcommand{\stepmark}[1]{\ifstepmark
%       \,\,\colorbox{Green!40!}{[{#1}\,point\ifdim#1pt=1pt%
%          \else
%             s
%          \fi%
%       ]}
%    \else%
%       {}
%    \fi%
% }

\newcommand{\blankpages}[1]{
   \ifthenelse{\boolean{printanswers}}
   {}
   {
      \foreach \n [count=\ni] in {1,...,#1}{%
         \ifnum\ni=#1% 
            \newpage%
         \else%
            \newpage%
            \mbox{}%
         \fi%
      }
   }
}

\noprintanswers
%\stepmarktrue
%\stepmarkfalse

\begin{document}

\textsc{{\normalsize\courseno \push \testname}\smallskip\\
{\normalsize\coursename \push \testtime}\smallskip\\
{\normalsize\instiname \push \testdate}}
\begin{center}
   \fbox{\large{\textsc{\textbf{PART-A}}}}
\end{center}

\ifthenelse{\boolean{printanswers}}
{
   \fullwidth{\hrulefill\raisebox{-2pt}{\fbox{\textsc{Model Answers of
               \testname\ (Total points: \numpoints)}}}\hrulefill}
}
{
   \hrulefill\\[5mm]
   Name of the Student: \hrulefill\\[5mm]
   Signature of the Student: \hrulefill \hspace{10mm} Roll No.: \hrulefill\\[5mm]
   Signature of the Invigilator: \hrulefill
   \fullwidth{\hrulefill\raisebox{-2pt}{\footnotesize\fbox{\textsc{Instructions}}}\hrulefill}
   \begin{enumerate}
      % \item To `print' means to write legibly in capital letters.
      % \item Put your signature in the space provided:$\rightarrow$ \push
      % \framebox{\rule{100pt}{0pt}\rule{0pt}{10pt}} 
      % \item Print your roll number:$\rightarrow$ \push \emptyrollno
      % \item Print your name as registered:\push \emptyname
      \item Check that you have \textbf{\numpages\  printed pages}. The PART-A
         of the examination has \textbf{\numquestions\ questions}, for a total of
         \textbf{\numpoints\ points}.
      \item Attempt \textbf{all questions}.
      \item There is \textbf{no credit} for a solution if the appropriate work is not
         shown, even if the answer is correct.
      \item Write your answers in this booklet and only in the \textbf{space marked for
         answers}. Any answer written in any space other than the designated space will
         \textbf{not be evaluated}.
      \item The pages numbered \textbf{8 to 10} are kept as \textbf{additional pages}. If
         you cannot able to fit a complete answer in the space provided just
         after the question, you can use these three pages by referring the page
         number of additional page in the original space.
      \item At the beginning of the examination, a supplementary sheet has been
         provided only for rough work. Should you need more, please ask the
         invigilator.
      \item No question requires any clarification from the instructors or
         invigilators, even if a question has an error or incomplete data. The
         students are advised to write answer according to their understanding
         or write reasons for why it is not possible to solve partially or
         completely that question by citing errors/insufficient data. 
   \end{enumerate}

   \fullwidth{\hrulefill\raisebox{-2pt}{\footnotesize\fbox{\textsc{For Office Use Only}}}\hrulefill}

   \vspace{0.2in}
   \begin{center}
      \cellwidth{12mm} 
      \gradetablestretch{3.0}
      \gradetable[h][questions]
   \end{center}

   \newpage

   \fullwidth{\hrulefill\raisebox{-2pt}{\footnotesize\fbox{\textsc{Questions}}}\hrulefill}
}

\begin{questions}

   \question[5] Let $Y$ and $\epsilon$ are the response variable and the error
   in a regression problem, respectively. $\theta_1$ and $\theta_2$ are the two
   parameters. Let $y_i, i=1,2,3$, be the observed values of the response
   variable. Consider,
   \begin{eqnarray}
      Y_1 &=& \theta_1 + \theta_2 + \epsilon_1, \nonumber \\
      Y_2 &=& \theta_1 - 2\theta_2 + \epsilon_2, \nonumber \\
      Y_3 &=& 2\theta_1 - \theta_2 + \epsilon_3, \nonumber 
   \end{eqnarray}
   where $E(\epsilon_i) = 0, \mbox{ and } V(\epsilon_i)= \sigma^2, i=1,2,3;
   \sigma > 0$. Find the least square estimates (LSE) of $\theta_1$ and
   $\theta_2$. 

   \begin{solution}
      We have
      \begin{align*}
         \begin{pmatrix} y_1 \\ y_2 \\ y_3 \end{pmatrix}_{3 \times
         1} = \begin{bmatrix} 1 & 1 \\ 1 & -2 \\ 2 & -1 \end{bmatrix}_{3 \times 2}
         \begin{bmatrix} \theta_1 \\ \theta_2 \end{bmatrix}_{2 \times 1} +
         \begin{pmatrix} \epsilon_1 \\ \epsilon_2 \\ \epsilon_3 \end{pmatrix} 
         \implies
         Y = X\theta + \epsilon.
      \end{align*}
      The least square estimator of $\theta$ is $ \hat{\theta} = (X'X)^{-1}X'Y$.
      Now,
      \begin{align*}
         (X'X) &= \begin{bmatrix} 1 & 1 & 2 \\ 1 & -2 & -1 \end{bmatrix}_{2
         \times 3} \begin{bmatrix} 1 & 1 \\ 1 & -2 \\ 2 & -1 \end{bmatrix}_{3
         \times 2} \\
         &= \begin{bmatrix} 1+1+4=6 & 1-2-2=-3 \\ 1-2-2=-3 & 1+4+1=6
         \end{bmatrix} = \begin{bmatrix} 6 & -3 \\ -3 & 6 \end{bmatrix}. \\
         (X'X)^{-1} &= \frac{1}{36 - 9} \begin{bmatrix} 6 & 3 \\ 3 & 6
         \end{bmatrix} = \frac{3}{27} \begin{bmatrix} 2 & 1 \\ 1 & 2
         \end{bmatrix} = \frac{1}{9} \begin{bmatrix} 2 & 1 \\ 1 & 2 \end{bmatrix}
      \end{align*}
      Thus
      \begin{align*}
         \hat{\theta} 
         &= \frac{1}{9} \begin{bmatrix} 2 & 1 \\ 1 & 2
         \end{bmatrix}_{2 \times 2} \begin{bmatrix} 1 & 1 & 2 \\ 1 & -2 & -1
         \end{bmatrix}_{2 \times 3} \begin{bmatrix} y_1 \\ y_2 \\ y_3
         \end{bmatrix} \\
         &= \frac{1}{9} \begin{bmatrix} 2 & 1 \\ 1 & 2 \end{bmatrix}
         \begin{bmatrix} y_1 + y_2 + 2y_3 \\ y_1 - 2y_2 - y_3 \end{bmatrix} \\
         &= \frac{1}{9} \begin{bmatrix} 2(y_1 + y_2 + 2y_3) + (y_1 - 2y_2 - y_3)
         \\ (y_1 + y_2 + 2y_3) + 2(y_1 - 2y_2 - y_3) \end{bmatrix} \\
         &= \frac{1}{9} \begin{bmatrix} 3y_1 + 3y_3 \\ 3y_1 - 3y_2
         \end{bmatrix}\\
         &= \begin{bmatrix} \frac{1}{3}(y_1 + y_3) \\ \frac{1}{3}(y_1 - y_2)
         \end{bmatrix}.
      \end{align*}

      Hence,
      $ \hat{\theta}_1 = \frac{1}{3}(y_1 + y_3) $
      and 
      $ \hat{\theta}_2 = \frac{1}{3}(y_1 - y_2) $.
   \end{solution}

   \blankpages{2}

   \question Let $U_1$ and $U_2$ be two independent uniform random variables on the
   interval $(0,\,1)$. Suppose that $\boldsymbol{X} = \left( X_1,\,X_2,\,X_3
   \right)^\prime$, where $X_1=U_1$, $X_2=U_2$, and $X_3=U_1+U_2$. 
   \begin{parts}
      \part[3] Compute the correlation matrix $\boldsymbol{\rho}$ of $\boldsymbol{X}$.

      \begin{solution}
      
         To compute correlation matrix $\rho$ of $X$.

         \begin{align*}
            E(U_1) = E(U_2) &= \frac{1}{2} && \text{ and } && V(U_1) = V(U_2) =
            \frac{1}{12}.
         \end{align*}
         So, $V(X_1) = V(X_2) = \frac{1}{12}$ and $V(X_3) = V(U_1) + V(U_2) =
         \frac{2}{12}$.
         \begin{align*}
            Cov(X_1, X_2) &= Cov(U_1, U_2) = 0 = Cov(X_2, X_1) \\ %\dots (1) \quad
            %(\because U_1 \text{ \& } U_2 \text{ are ind.}) \\
            Cov(X_1, X_3) &= Cov(U_1, U_1 + U_2) \\
                          &= Cov(U_1, U_1) \\
                          &= V(U_1) = \frac{1}{12} = Cov(X_3, X_1) \\
            Cov(X_2, X_3) &= Cov(U_2, U_1 + U_2) \\
                          &= V(U_2) = \frac{1}{12} = Cov(X_3, X_2).
         \end{align*}
         Now, $Corr(X_i, X_i) = 1$ for $i=1,2,3$. Moreover,
         \begin{align*}
            Corr(X_1, X_2) &= 0, \\
            Corr(X_1, X_3) &= Corr(X_3, X_1) = \frac{Cov(X_1, X_3)}{\sqrt{V(X_1) V(X_3)}} =
                           \frac{\frac{1}{12}}{\sqrt{\frac{1}{12} \cdot
                           \frac{2}{12}}} = \frac{1}{\sqrt{2}}, \\
            Corr(X_2, X_3) &= \frac{Cov(X_2, X_3)}{\sqrt{V(X_2) V(X_3)}} =
            \frac{1}{\sqrt{2}}.
         \end{align*}

         Hence, the correlation matrix $\rho$ of $X$ is 
         $$ \rho = \begin{bmatrix} 1 & 0 & \frac{1}{\sqrt{2}} \\ 0 & 1 &
         \frac{1}{\sqrt{2}} \\ \frac{1}{\sqrt{2}} & \frac{1}{\sqrt{2}} & 1
         \end{bmatrix} .$$

      \end{solution}

      \part[3] Find the variance of first principal component based on the
      correlation matrix that you find in part (a).

      \begin{solution}
         Let $Y_1$ be the first principal component based on the correlation matrix $\rho$. Then
         $$ V(Y_1) = \lambda_1,$$
         where $\lambda_1$ is the highest eigenvalue of the matrix $\rho$.  The
         eigenvalues of $\rho$ are solution of the following equation:
         \begin{align*}
            & \quad (1-\lambda) \begin{vmatrix} 1-\lambda & \frac{1}{\sqrt{2}} \\
            \frac{1}{\sqrt{2}} & 1-\lambda \end{vmatrix} - 0 \begin{vmatrix} 0 &
            \frac{1}{\sqrt{2}} \\ \frac{1}{\sqrt{2}} & 1-\lambda \end{vmatrix} +
            \frac{1}{\sqrt{2}} \begin{vmatrix} 0 & 1-\lambda \\
            \frac{1}{\sqrt{2}} & \frac{1}{\sqrt{2}} \end{vmatrix} = 0\\
            \Rightarrow & \quad (1-\lambda) \left[ (1-\lambda)^2 - \frac{1}{2} \right] -
            0 + \frac{1}{\sqrt{2}} \left[ 0 - \frac{1}{\sqrt{2}}(1-\lambda)
            \right] = 0 \\
            \Rightarrow & \quad (1-\lambda)^3 - \frac{1}{2}(1-\lambda) - \frac{1}{2}(1-\lambda) = 0 \\
            \Rightarrow & \quad (1-\lambda)^3 - (1-\lambda) = 0 \\
            \Rightarrow & \quad (1-\lambda) \left( (1-\lambda)^2 - 1 \right) = 0 \\
            \Rightarrow & \quad (1-\lambda) (\lambda^2 - 2\lambda) = 0 \\
            \Rightarrow & \quad \lambda (1-\lambda) (\lambda - 2) = 0.
         \end{align*}
         Hence, the largest eigenvalue is $2$. So,
         $$ V(Y_1) = 2. $$
      \end{solution}

   \end{parts}

   \blankpages{2}

   \question
   Let the matrix of distance be
   \[
      \begin{blockarray}{ccccc}
         & 1 & 2 & 3 & 4  \\
         \begin{block}{c[cccc]}
            \addlinespace
            1 & 0 &   &   &   \\
            \addlinespace
            2 & 1 & 0 &   &   \\
            \addlinespace
            3 & 11& 2 & 0 &   \\
            \addlinespace
            4 & 5 & 3 & 3.5 & 0  \\
            \addlinespace
         \end{block}
      \end{blockarray}
   \]

   \begin{parts}
      \part[6] Cluster the four items using each of the following procedures:
      \begin{subparts}
         \subpart Complete linkage hierarchical procedure.
         \subpart Average linkage hierarchical procedure.
      \end{subparts}

      \begin{solution}
         %          \[
         %             D = 
         %             \begin{array}{cc}
         %  & \begin{array}{cccc} 1 & 2 & 3 & 4 \end{array} \\
         %                \begin{array}{c} 1 \\ 2 \\ 3 \\ 4 \end{array} &
         %                \left[ \begin{array}{cccc}
         %                      0 & & & \\
         %                      1 & 0 & & \\
         %                      11 & 2 & 0 & \\
         %                      5 & 3 & 4 & 0
         %                \end{array} \right]
         %             \end{array}
         %             = (d_{ij}) \quad i,j=1,2,3,4.
         %          \]
         %
         %          \vspace{0.5cm}
         %          \noindent \textbf{(a) Using Single linkage,}
         %          \[ \min(d_{ij}) = 1 = d_{21} \]
         %          So, we merge objects 1 \& 2 to form cluster (12). Remaining objects 3, 4. \\
         %          The nearest neighbor distances are
         %          \begin{align*}
         %             d_{(12),3} &= \min(d_{13}, d_{23}) = \min(11, 2) = 2. \\
         %             d_{(12),4} &= \min(d_{14}, d_{24}) = \min(5, 3) = 3.
         %          \end{align*}
         %          So we obtain the new distance matrix,
         %          \[
         %             \begin{array}{cc}
         %  & \begin{array}{ccc} (12) & 3 & 4 \end{array} \\
         %                \begin{array}{c} (12) \\ 3 \\ 4 \end{array} &
         %                \left[ \begin{array}{ccc}
         %                      0 & & \\
         %                      2 & 0 & \\
         %                      3 & 4 & 0
         %                \end{array} \right]
         %             \end{array}
         %          \]
         %          \[ \min(d_{ij}) = d_{(12),3} = 2. \]
         %          So, we merge cluster (3) with cluster (12). \\
         %          Now the nearest-neighbor distances are
         %          \begin{align*}
         %             d_{(123),4} &= \min(d_{(12),4}, d_{34}) \\
         %                         &= \min(3, 4) \\
         %                         &= 3
         %          \end{align*}
         %          So, we obtain the new distance matrix,
         %          \[
         %             \begin{array}{cc}
         %  & \begin{array}{cc} (123) & 4 \end{array} \\
         %                \begin{array}{c} (123) \\ 4 \end{array} &
         %                \left[ \begin{array}{cc}
         %                      0 & \\
         %                      3 & 0
         %                \end{array} \right]
         %             \end{array}
         %          \]
         %          Now cluster (123) \& (4) are merged together to form a single cluster.
         %
         %          \vspace{0.5cm}
         %          \noindent \textbf{(b) Using Complete linkage,}
         % \[
         %    D = 
         %    \begin{array}{cc}
         %       & \begin{array}{cccc} 1 & 2 & 3 & 4 \end{array} \\
         %       \begin{array}{c} 1 \\ 2 \\ 3 \\ 4 \end{array} &
         %       \left[ \begin{array}{cccc}
         %             0 & & & \\
         %             1 & 0 & & \\
         %             11 & 2 & 0 & \\
         %             5 & 3 & 4 & 0
         %       \end{array} \right]
         %    \end{array}
         % \]
         At the first stage merge objects 1 and 2 together as they are most similar. \\
         At stage 2, we compute distances using complete linkage:
         \begin{align*}
            d_{(12),3} &= \max(d_{13}, d_{23}) = \max(11, 2) = 11,\\
            d_{(12),4} &= \max(d_{14}, d_{24}) = \max(5, 3) = 5.
         \end{align*}
         So, we obtain the distance matrix,
         \[
            \begin{array}{cc}
               & \begin{array}{ccc} (12) & 3 & 4 \end{array} \\
               \begin{array}{c} (12) \\ 3 \\ 4 \end{array} &
               \left[ \begin{array}{ccc}
                     0 & & \\
                     11 & 0 & \\
                     5 & 3.5 & 0
               \end{array} \right]
            \end{array}
         \]
         At stage 3, merge objects 3 \& 4, as they are most similar. \\
         Now, we compute,
         \begin{align*}
            d_{(12)(34)} &= \max\{d_{(12)3}, d_{(12)4}\} \\
                          &= \max\{11, 5\} \\
                          &= 11
         \end{align*}
         So, we obtain the new distance matrix,
         \[
            \begin{array}{cc}
            & \begin{array}{cc} (12) & (34) \end{array} \\
               \begin{array}{c} (12) \\ (34) \end{array} &
               \left[ \begin{array}{cc}
                     0 & \\
                     11 & 0
               \end{array} \right]
            \end{array}
         \]
         Now, at the final stage, merge clusters (12) \& (34) together to form a
         single cluster.

         % \vspace{0.5cm}
         % \noindent \textbf{(c) Dendogram from (a)}
         %
         % \vspace{0.3cm}
         % \begin{tikzpicture}[yscale=1.5]
         %    % Y-axis
         %    \draw (0,0) -- (0,3.5); 
         %    \foreach \y in {1,2,3} \draw (-0.1,\y) -- (0,\y) node[left] {\y};
         %
         %    % Nodes
         %    \node[below] at (1,0) {1};
         %    \node[below] at (2,0) {2};
         %    \node[below] at (3,0) {3};
         %    \node[below] at (4,0) {4};
         %
         %    % Links
         %    \draw (1,0) -- (1,1) -- (2,1) -- (2,0);
         %    \draw (1.5,1) -- (1.5,2) -- (3,2) -- (3,0);
         %    \draw (2.25,2) -- (2.25,3) -- (4,3) -- (4,0);
         %
         %    % Points
         %    \fill (1,0) circle (1.5pt);
         %    \fill (2,0) circle (1.5pt);
         %    \fill (3,0) circle (1.5pt);
         %    \fill (4,0) circle (1.5pt);
         %    \fill (1.5,1) circle (1.0pt);
         %    \fill (2.25,2) circle (1.0pt);
         % \end{tikzpicture}
         %
         % \vspace{0.5cm}
         % \noindent \textbf{Dendogram from (b)}
         %
         % \vspace{0.3cm}
         
         \textbf{jaa}

         At the first stage merge objects 1 and 2 together as they are most similar. \\
         At stage 2, we compute distances using complete linkage:
         \begin{align*}
            d_{(12),3} &= \frac{d_{13} + d_{23}}{2} = \frac{11 + 2}{2} = 6.5,\\
            d_{(12),4} &= \frac{d_{14} + d_{24}}{2} = \frac{5 + 3}{2} = 4.
         \end{align*}
         So, we obtain the distance matrix,
         \[
            \begin{array}{cc}
               & \begin{array}{ccc} (12) & 3 & 4 \end{array} \\
               \begin{array}{c} (12) \\ 3 \\ 4 \end{array} &
               \left[ \begin{array}{ccc}
                     0 & & \\
                     6.5 & 0 & \\
                     4 & 3.5 & 0
               \end{array} \right]
            \end{array}
         \]
         At stage 3, merge objects 3 \& 4, as they are most similar. \\
         Now, we compute,
         \begin{align*}
            d_{(12)(34)} &= \frac{d_{(12)3} + d_{(12)4}}{2} = \frac{6.5 + 4}{2}
            = 5.25.
         \end{align*}
         So, we obtain the new distance matrix,
         \[
            \begin{array}{cc}
               & \begin{array}{cc} (12) & (34) \end{array} \\
               \begin{array}{c} (12) \\ (34) \end{array} &
               \left[ \begin{array}{cc}
                     0 & \\
                     5.25 & 0
               \end{array} \right]
            \end{array}
         \]
         Now, at the final stage, merge clusters (12) \& (34) together to form a
         single cluster.
      \end{solution}

      \part[3] Draw the dendrograms and compare the results obtained using previous
      two procedures.

      \begin{solution}
         Dendogram for complete linkage:
         \begin{center}
         \begin{tikzpicture}[yscale=0.5]
            % Y-axis
            \draw (0,0) -- (0,12); 
            \foreach \y in {1,3.5,11} \draw (-0.1,\y) -- (0,\y) node[left] {\y};

            % Nodes
            \node[below] at (1,0) {1};
            \node[below] at (2,0) {2};
            \node[below] at (3,0) {3};
            \node[below] at (4,0) {4};

            % Links
            \draw (1,0) -- (1,1) -- (2,1) -- (2,0);
            \draw (3,0) -- (3,4) -- (4,4) -- (4,0);
            \draw (1.5,1) -- (1.5,11) -- (3.5,11) -- (3.5,4);

            % Points
            \fill (1,0) circle (1.5pt);
            \fill (2,0) circle (1.5pt);
            \fill (3,0) circle (1.5pt);
            \fill (4,0) circle (1.5pt);
            \fill (1.5,1) circle (1.0pt);
            \fill (3.5,4) circle (1.0pt);
            \fill (2.5,11) circle (1.0pt);
         \end{tikzpicture}
         \end{center}
         Dendogram for average linkage:
         \begin{center}
         \begin{tikzpicture}%[yscale=0.3]
            % Y-axis
            \draw (0,0) -- (0,6); 
            \foreach \y in {1,3.5,5.25} \draw (-0.1,\y) -- (0,\y) node[left] {\y};

            % Nodes
            \node[below] at (1,0) {1};
            \node[below] at (2,0) {2};
            \node[below] at (3,0) {3};
            \node[below] at (4,0) {4};

            % Links
            \draw (1,0) -- (1,1) -- (2,1) -- (2,0);
            \draw (3,0) -- (3,4) -- (4,4) -- (4,0);
            \draw (1.5,1) -- (1.5,5.25) -- (3.5,5.25) -- (3.5,4);

            % Points
            \fill (1,0) circle (1.5pt);
            \fill (2,0) circle (1.5pt);
            \fill (3,0) circle (1.5pt);
            \fill (4,0) circle (1.5pt);
            \fill (1.5,1) circle (1.0pt);
            \fill (3.5,4) circle (1.0pt);
            \fill (2.5,5.25) circle (1.0pt);
         \end{tikzpicture}
         \end{center}

         \vspace{0.5cm}
         % \noindent \textbf{Results comparing both (a) \& (b).} \\
         Comparing the two dendrograms, we observe that the single linkage and
         complete linkage do not differ in allocating objects in this case.
      \end{solution}
   \end{parts}

   \blankpages{2}

   \ifthenelse{\boolean{printanswers}}
   {
   }
   {
      \fullwidth{\hrulefill\raisebox{-2pt}{\footnotesize\fbox{\textsc{Additional
      Pages for Answers}}}\hrulefill}
   }

   \blankpages{3}

\end{questions}

\end{document}
